\section*{3. Complementos de álgebra homológica}

El objeto principal de esta sección es definir el isomorfismo
\[
 \Biext^1(P,Q;G)\simeq
 \Ext^1(P\overset{\mathbb L}{\otimes}Q,G)
\]
anunciado anteriormente. Para ello necesitaremos ciertos desarrollos
de álgebra homológica que tienen interés independiente y ocuparán las
secciones 3.1 a 3.4. Los desarrollos de 3.1 y 3.3 pueden considerarse
complementos de [4], especialmente de su sección 6.9.

\subsection*{3.1. Construcción de una categoría cofibrada aditiva
$\underline\Psi_{L_\bullet}$ a partir de un complejo de cadenas
$L_\bullet$ en una categoría abeliana $\underline C$}

Sea $\underline C$ una categoría abeliana y sea $L_\bullet$ un
complejo de cadenas en $\underline C$,
\begin{equation}\tag{3.1.1}
 L_2\longrightarrow L_1\longrightarrow L_0\longrightarrow0,
\end{equation}
escrito con la convención de grados inferiores, de modo que el
diferencial tiene grado $-1$. Así, $L_i=0$ para $i<0$. Definiremos una
categoría cofibrada
\[
 \underline\Psi=\underline\Psi_{L_\bullet}
\]
sobre $\underline C$; salvo equivalencia canónica de categorías
cofibradas, es la categoría del mismo nombre de [4, 6.9]. Los objetos
de $\underline\Psi_{L_\bullet}$ son los pares
\[
 (X,\alpha),
\]
donde $X$ es una extensión de $L_0$ por un objeto $A$ de
$\underline C$, que depende de $X$,
\begin{equation}\tag{3.1.2}
 0\longrightarrow A\xrightarrow{i}X\xrightarrow{j}L_0
 \longrightarrow0,
\end{equation}
y $\alpha$ es una trivialización de la imagen inversa $d_0^*(X)$ de
esta extensión por $d_0:L_1\longrightarrow L_0$. Se exige que la
trivialización correspondiente de la imagen inversa
$d_1^*d_0^*(X)$ por $d_1:L_2\longrightarrow L_1$ sea la evidente,
obtenida del isomorfismo de transitividad
\[
 d_1^*d_0^*(X)\simeq(d_0d_1)^*(X)
\]
y de la relación $d_0d_1=0$.

Equivalentemente, $\alpha$ es un levantamiento de
$d_0:L_1\longrightarrow L_0$ a un morfismo
$\alpha:L_1\longrightarrow X$, sujeto a $\alpha d_1=0$:
\begin{equation}\tag{3.1.3}
\begin{tikzcd}[row sep=large,column sep=large]
 &&0\arrow[d]\\
 &&A\arrow[d,"i"]\\
 &&X\arrow[d,"j"]\\
 L_2\arrow[r,"d_1"']&L_1\arrow[r,"d_0"']\arrow[ru,"\alpha"]
   &L_0\arrow[d]\\
 &&0
\end{tikzcd}
\qquad,\qquad \alpha d_1=0.\quad(*)
\end{equation}
He aquí una descripción más sencilla, señalada por Deligne:
$\underline\Psi_{L_\bullet}(A)$ es la categoría de las extensiones del
complejo $L_\bullet$ por $A$, considerado como un complejo concentrado
en grado cero.

En esta forma se ve que la descripción de los objetos $(X,\alpha)$
depende solo del complejo de cadenas $[L_\bullet]$
\underline{truncado en grado 1}, obtenido de $L_\bullet$ ``matando
$H_i(L_\bullet)$ para $i\geq2$'': sus términos son cero en grados al
menos $2$, coinciden con los de $L_\bullet$ en grados a lo sumo $0$, y
su término de grado $1$ es el conúcleo de
$d_1:L_2\longrightarrow L_1$.

Un homomorfismo de $(X,\alpha)$ a $(X',\alpha')$ es un homomorfismo de
extensiones $u:X\longrightarrow X'$, compatible con la identidad
sobre $L_0$, tal que $u\alpha=\alpha'$. Los homomorfismos se componen
componiendo los homomorfismos de extensiones correspondientes. Esto da
la categoría anunciada $\underline\Psi_{L_\bullet}$. Salvo
isomorfismo canónico, depende únicamente del truncamiento en grado
$1$, $[L_\bullet]$, de $L_\bullet$.

Hay un funtor evidente $(X,\alpha)\longmapsto A$,
\begin{equation}\tag{3.1.4}
 \pi:\underline\Psi_{L_\bullet}\longrightarrow\underline C.
\end{equation}
\underline{Este funtor es cofibrado}. En efecto, sea
$u:A\longrightarrow B$ un morfismo de $\underline C$, y sea
$(X,\alpha)$ un objeto de la fibra
$\underline\Psi_{L_\bullet}(A)$. Así, $X$ es una extensión de $L_0$
por $A$. Formemos la extensión $u_*(X)$ de $L_0$ por $B$, imagen
directa de $X$ por $u$, y dotémosla de la trivialización parcial
$\beta=f\alpha$, donde $f:X\longrightarrow u_*(X)$ es el morfismo
canónico. Se sigue inmediatamente que
\[
 f:(X,\alpha)\longrightarrow(u_*(X),\beta)
\]
es cocartesiano para (3.1.4): posee la propiedad universal usual para
los $u$-morfismos de $\underline\Psi_{L_\bullet}$ con fuente
$(X,\alpha)$. Esta no es sino la propiedad universal correspondiente,
bien conocida y elemental, de $f:X\longrightarrow u_*(X)$.

Finalmente, la categoría cofibrada $\underline\Psi_{L_\bullet}$ es
\underline{aditiva}; es decir, satisface
\begin{equation}\tag{3.1.5}
 \underline\Psi_{L_\bullet}(0)
 \quad\text{es equivalente a la categoría con un solo objeto},
\end{equation}
y, para $A,B$ arbitrarios en $\underline C$,
\begin{equation}\tag{3.1.6}
 \underline\Psi_{L_\bullet}(A\times B)
 \longrightarrow
 \underline\Psi_{L_\bullet}(A)\times
 \underline\Psi_{L_\bullet}(B)
 \quad\text{es una equivalencia}.
\end{equation}
Esto se deduce inmediatamente de las propiedades correspondientes de
la categoría cofibrada de las extensiones de $L_0$ por objetos
variables de $\underline C$. Utilizaremos libremente las consecuencias
de la aditividad de $\underline\Psi_{L_\bullet}$ establecidas en
[4, 1] y recordadas anteriormente en 2.5.

3.1.7. Entre estas consecuencias, para todo objeto $A$ de
$\underline C$, la categoría $\underline\Psi_{L_\bullet}(A)$ tiene un
objeto \underline{cero}, o \underline{trivial} (2.5.1). Es la
extensión trivial $X=L_0+A$ de $L_0$ por $A$, provista de la
trivialización parcial $\alpha:L_1\longrightarrow X$ cuyas componentes
son $d_0$ y $0$.

Además, por 2.5.2, para todo $(X,\alpha)$ de
$\underline\Psi_{L_\bullet}(A)$, el grupo de automorfismos de este
objeto es canónicamente isomorfo al grupo de automorfismos del objeto
trivial. Denotemos este último grupo conmutativo por
$\Psi^0_{L_\bullet}(A)$. La descripción explícita anterior da los
isomorfismos canónicos
\begin{multline}\tag{3.1.8}
 \Psi^0_{L_\bullet}(A)
 =\Ker\bigl(\Hom(L_0,A)\longrightarrow\Hom(L_1,A)\bigr)
 \simeq\Hom(H_0(L_\bullet),A)\\
 \simeq\Ext^0(L_\bullet,A),
\end{multline}
inducidos por la identificación usual del grupo de automorfismos de
la extensión $X$ con $\Hom(L_0,A)$. El isomorfismo
\sisi{(3.1.7.1)}{(3.1.8)} es manifiestamente funtorial en $A$.

3.1.9. Finalmente, recordemos de 2.5.2 que la estructura
multiplicativa natural sobre la categoría fibra
$\underline\Psi_{L_\bullet}(A)$ dota a su conjunto de clases de
isomorfismo de una estructura de grupo natural. Denotamos este grupo
por $\Psi^1_{L_\bullet}(A)$; es un funtor aditivo de $A$.
